% TOP_PROBLEM: {"id": 30006935, "title": "Von Neumann's Problem for II_1 Factors", "category_id": 8, "category": {"id": 8, "name": "analysis", "display_name": "Analysis", "description": "Limits, continuity, calculus, and function theory.", "slug": "analysis", "order_index": 8, "created_at": "2026-07-31T15:26:25.671Z"}, "status": "open"}
% FIRST_POSTED: 2026-09-27
% ATTEMPT_STATUS: unresolved
% !TeX program = lualatex
\documentclass[11pt]{article}
\usepackage[margin=25mm]{geometry}
\usepackage{fontspec}
\setmainfont{Latin Modern Roman}
\usepackage{amsmath,amssymb,mathtools}
\usepackage{unicode-math}
\setmathfont{Latin Modern Math}
\usepackage{xurl,hyperref}
\hypersetup{colorlinks=true,urlcolor=blue}
\setcounter{secnumdepth}{0}
\setlength{\parindent}{0pt}
\setlength{\parskip}{5pt}
\setlength{\emergencystretch}{3em}
\begin{document}
\section*{\texorpdfstring{Von Neumann's Problem for II\_1 Factors}{Von Neumann's Problem for II\_1 Factors}}\label{30006935}
\subsection{Definitions and mathematical statement}
A $\mathrm{II}_1$ factor is a von Neumann algebra with center ${\mathbb{C}}1$, a faithful normal tracial state, and no nonzero minimal projections. Here amenability is the usual injectivity/amenability property of von Neumann algebras. For a group $G$, $L(G)$ is the weak operator closure of the algebra generated by its left translation operators on $\ell^2(G)$.

Let $F_2$ be the free group on two generators. The question is whether every nonamenable $\mathrm{II}_1$ factor $M$ contains a unital von Neumann subalgebra isomorphic to $L(F_2)$, or whether a counterexample exists. Such a subalgebra is itself a factor, so this is also a subfactor-containment question. No finite-index hypothesis is added. The catalogue imposes no separability assumption, and none is silently supplied here.
\par\smallskip\textit{Formulation and concept references:} \href{https://math.vanderbilt.edu/peters10/problems.html}{[S1]}.

\subsection{Short English statement}
Does every nonamenable finite factor contain the von Neumann algebra of the free group on two generators?
\subsection{Research attempt}
\paragraph{Scope, attribution, and source audit.}
This unresolved attempt was prepared by GPT-6 Astra Ultra on 27 September
2026 through primary-source review, explicit group-word constructions,
and checks of convergence and trace conventions. The question in
Peterson's list [S1] concerns actual unital containment in the given
factor, with no separability or index restriction. It is stronger than
containment in an ultrapower. The attempt reduces containment to an
exact moment problem and constructs finite moment matches inside
amenable matrix algebras. That construction is a test of an attempted
proof, not a counterexample to the problem, since its ambient examples
are amenable.

\paragraph{An exact certificate using two unitaries.}
Let $(M,\tau)$ be any finite von Neumann algebra with faithful normal
trace. A sufficient and necessary certificate for a trace-preserving
copy of $L(F_2)$ is a pair of unitaries $u,v\in M$ such that
\begin{equation}\label{a439:words}
 \tau(w(u,v))=0
 \quad\text{for every nonidentity reduced word }
 w\in F_2.
\end{equation}
The canonical generators of $L(F_2)$ satisfy this because the
canonical trace extracts the identity coefficient. Conversely define
$\pi(a)=u$, $\pi(b)=v$ on the free group. For group-algebra elements,
\eqref{a439:words} gives
\[
 \tau(\pi(x)^*\pi(y))=\langle x,y\rangle_{\ell^2(F_2)}.
\]
Thus the map sending $\delta_g$ to $\pi(g)$ extends to a unitary
from $\ell^2(F_2)$ onto $L^2(N,\tau)$, where $N=W^*(u,v)$.
It intertwines left multiplication by every group generator
with the corresponding operator on $L^2(N)$. Taking the von
Neumann closures therefore identifies $L(F_2)$ normally and
trace-preservingly with $N$. Faithfulness of the traces makes
this identification injective, and the common identity makes
it unital.

In a factor a unital copy automatically has the correct normalized
trace, because the trace restricted to a $\mathrm{II}_1$ factor
is its unique normalized normal trace. Hence no extra choice of
trace weakens the requested containment. Equivalently $u,v$ are
free Haar unitaries: the reduced-word condition includes all
nonzero powers of each generator and all alternating products
of their centered powers.

\paragraph{A complete restricted class: group factors with a free subgroup.}
If a group $G$ contains elements $a,b$ generating a subgroup
isomorphic to $F_2$, its canonical translation unitaries in
$L(G)$ satisfy \eqref{a439:words}. Every nonidentity reduced
word is a nonidentity group element and therefore has trace
zero. This proves containment for any such group factor,
with no finite-index condition. If $L(G)$ is a
$\mathrm{II}_1$ factor, it is a genuine instance of the question.

The contrapositive is unavailable. A nonamenable group without
a free subgroup does not automatically yield a counterexample:
unitaries generating a copy of $L(F_2)$ inside $L(G)$ need
not be individual group elements. They may be operators
with large Fourier support. A group-theoretic obstruction
to selecting two group generators does not exclude the
von Neumann subalgebra certificate above.

\paragraph{Finite moment tests can all hold in finite algebras.}
Here is an explicit residual-finiteness construction requiring
no asymptotic random-matrix theorem. Fix $N$. Let $B_N$ be
the finite set of reduced words of length at most $N$.
Left multiplication by $a$ defines a partial bijection of
$B_N$ wherever both the input and reduced output remain in
$B_N$. Complete it to a permutation $A$ of $B_N$: the
unused domain and range have the same cardinality.
Do the same for $b$, obtaining $B$.

Every nonidentity reduced word $w$ of length at most $N$
acts nontrivially on the empty word under these permutations.
Indeed the successive reduced suffixes obtained while
evaluating $w(A,B)$ have length at most $N$, so each
prescribed partial multiplication agrees with the free
group action, including inverse steps. The final image
is $w$, not the empty word. Thus in the finite permutation
group $Q_N=\langle A,B\rangle$ none of these words is
the identity.

Now take the left regular representation of $Q_N$ on
$\mathbb C^{|Q_N|}$. A nonidentity element permutes the
standard basis without fixed points, so its normalized
matrix trace is zero. Consequently the resulting unitaries
$U_N,V_N$ satisfy \eqref{a439:words} for every word of
length at most $N$, exactly, not just approximately.

They cannot satisfy all word conditions simultaneously
at this fixed matrix size. Such conditions would make
the infinitely many distinct group words orthonormal in
the finite-dimensional Hilbert space of matrices. The
finite-dimensionality obstruction is immediate even
without a polynomial identity for matrices.

\paragraph{Why passing to limits is the actual issue.}
The hyperfinite $\mathrm{II}_1$ factor $R$ contains unital
matrix algebras of every finite size, so the preceding
finite certificates can all be realized in $R$, one
pair for each $N$. Yet $R$ cannot contain $L(F_2)$:
amenability passes to expected subalgebras, every unital
von Neumann subalgebra of a finite tracial algebra has
a trace-preserving conditional expectation, and the
free group factor is nonamenable. These established
operator-algebra facts are recalled in [E1]. They are
used only to diagnose this proof strategy.

One can also see where an ultraproduct changes the
statement. The sequences $(U_N),(V_N)$ become a pair
with all required moments in a tracial matrix
ultraproduct, since each fixed word eventually has
trace zero. The ultraproduct supplies a compactness
device outside any particular original matrix
algebra or factor. It has not supplied a pair
inside the specified $M$.

Weak compactness inside $M$ does not fix the problem.
Weak limits of unitaries may fail to be unitary.
For example multiplication by $z^k$ in
$L^\infty(S^1)$ converges weak-star to zero, as
follows first for trigonometric polynomials and
then for $L^1$ tests by approximation. Products
are not jointly weakly continuous. Therefore
taking separate weakly convergent subsequences
and substituting their limits into every word
is not a valid moment argument.

\paragraph{A precise sufficient coherence condition.}
Suppose one can find pairs $(u_j,v_j)$ of unitaries
inside the same $M$ such that
\[
 \max_{|w|\leq j,\ w\ne e}|\tau(w(u_j,v_j))|\leq2^{-j},
 \qquad
 \sum_j\bigl(\|u_{j+1}-u_j\|_2+
                   \|v_{j+1}-v_j\|_2\bigr)<\infty.
\]
Then the pairs converge in $L^2$ to unitaries $u,v\in M$.
For clarity, a uniformly operator-norm bounded
$L^2$-Cauchy sequence has a limit in the same bounded
ball of $M$, using weak compactness and uniqueness
of the $L^2$ limit. Multiplication is continuous
in $L^2$ on bounded sets, so $u^*u=uu^*=1$ and
the same holds for $v$.

A telescoping expansion of a word of length $\ell$
bounds its $L^2$ difference by
\[
 \|w(u_j,v_j)-w(u,v)\|_2
 \leq\ell\max(\|u_j-u\|_2,\|v_j-v\|_2),
\]
with the same estimate for inverse letters because
adjoints preserve the $2$-norm. Since
$|\tau(x)|\leq\|x\|_2$, all moments pass to zero.
The exact certificate proves containment.

This lemma makes the first gap concrete: nonamenability
would have to supply an exact free pair, or an
approximation scheme with sufficient coherence in
the original algebra. Arbitrarily accurate finite
moment matching, even exact matching, is inadequate.
No implication from nonamenability to the displayed
summable perturbation condition is proved here.

\paragraph{Verification and outcome.}
The finite-word construction was checked by enumerating
short reduced words and their partial permutation paths.
The GNS argument keeps the full group law, not merely
pairwise orthogonality of $u$ and $v$. All constructions
use only countably many operators, but this does not
silently assume the ambient factor is separable or prove
a reduction from every nonseparable factor to a special
separable class.
\textbf{UNRESOLVED.} Exact moment and coherent-limit
criteria are proved, as is a restricted group-factor
case. The universal nonamenable-factor assertion and
the existence of a counterexample both remain open
in this attempt.

\subsection{Sources}
\begin{itemize}
\item[S1]Open problems in operator algebras. Accessed 2026-09-01.
\url{https://math.vanderbilt.edu/peters10/problems.html}.
\textit{Formulation source.}
\item[E1]Jesse Peterson, \emph{Notes on von Neumann algebras}, author course notes, especially finite traces, conditional expectations and amenability.
\url{https://math.vanderbilt.edu/peters10/teaching/spring2013/vonNeumannAlgebras.pdf}.
\textit{Author reference for the established operator-algebra inputs; consulted 27 September 2026. The finite-word construction and coherent-limit lemma are proved in the attempt.}
\end{itemize}
\end{document}
